\section*{ID9517: Gravitation: Spacetime Curvature and Lensing Anomalies: R\_μν - tfrac12R,g\_μν}
\paragraph{Metadata.}
PiID: 0330276347870 | RTEID: RTE:P36956.6245:T4.3445 | UUID: 8da3ae60-40dd-5c79-980f-72e8e3560615

\paragraph{Canonical Statement.}
In general relativity, gravity is the manifestation of spacetime curvature. Einstein’s field equations relate the Ricci curvature tensor \$R\_\{\textbackslash{}mu\textbackslash{}nu\}\$ to stress-energy as \$R\_\{\textbackslash{}mu\textbackslash{}nu\} - \textbackslash{}tfrac\{1\}\{2\}R,g\_\{\textbackslash{}mu\textbackslash{}nu\} = \textbackslash{}frac\{8\textbackslash{}pi G\}\{c\textasciicircum{}4\}T\_\{\textbackslash{}mu\textbackslash{}nu\}\$. Solutions to these equations (such as the Schwarzschild metric) predict how mass curves space and deflects light. A natural place to introduce a dimensionless ratio is by comparing observed gravitational effects to predicted ones from visible matter alone. For example, the deflection angle of light by a mass \$M\$ (in the weak-field limit) is approximately

\paragraph{Canonical Equation.}
\[
R_{\mu\nu} - \tfrac{1}{2}R,g_{\mu\nu} = \frac{8\pi G}{c^4}T_{\mu\nu}
\]

\paragraph{Dictionary.}
In general relativity, gravity is the manifestation of spacetime curvature.

\paragraph{Encyclopedia.}
Gravitation: Spacetime Curvature and Lensing Anomalies: R\_μν - tfrac12R,g\_μν is an algebraic definition, written R\_μν - tfrac12R,g\_μν = (8π G)/(c\textasciicircum{}4)T\_μν. It is an algebraic definition expressing the quantity directly in terms of the variables on its right-hand side. It is built from π, μ, ν, G, c, T\_, defining R\_μν - tfrac12R,g\_μν. Within the Trawin Zero Point registry it serves as one element of the framework's mathematical narrative.
